\documentclass[aspectratio=169,10pt]{beamer}
\makeatletter\def\input@path{{../../estilo/}{../estilo/}{estilo/}}\makeatother
\usepackage{devmobile}
\modulo[Flutter]{02}{Dart Essencial}{Variáveis e tipos, null safety, controle de fluxo, funções, coleções, classes, mixins, records, padrões e sealed classes}

\begin{document}
\capa

\begin{frame}{Objetivos do módulo}
  \begin{conceito}[Ao final deste módulo você será capaz de]
    \begin{itemize}
      \item Declarar variáveis com \cod{var}, \cod{final}, \cod{const} e \cod{late} e usar os tipos básicos.
      \item Escrever código seguro com \emph{null safety} (\cod{?}, \cod{??}, \cod{?.}, \cod{!}).
      \item Usar \cod{if}, \cod{switch} como expressão, laços e funções com parâmetros nomeados.
      \item Manipular \cod{List}, \cod{Set} e \cod{Map} com operações funcionais.
      \item Criar classes, construtores, interfaces, \emph{mixins} e extensões.
      \item Modelar dados com \emph{records}, \emph{patterns}, \cod{enum} e \cod{sealed class}.
    \end{itemize}
  \end{conceito}
  \begin{dica}[Onde praticar]
    \url{https://dartpad.dev} roda Dart (e Flutter) no navegador, sem instalar nada.
  \end{dica}
\end{frame}

\sumario

% =============================================================================
%  Fundamentos compartilhados --- Módulo 02: linguagens de programação
%  Incluído pelas trilhas Kotlin e Flutter com % =============================================================================
%  Fundamentos compartilhados --- Módulo 02: linguagens de programação
%  Incluído pelas trilhas Kotlin e Flutter com % =============================================================================
%  Fundamentos compartilhados --- Módulo 02: linguagens de programação
%  Incluído pelas trilhas Kotlin e Flutter com \input{../fundamentos/f02-linguagens}
% =============================================================================
\section{Fundamentos: linguagens de programação}

\begin{frame}{Paradigmas de programação}
  \begin{columns}[T]
    \begin{column}{0.54\textwidth}
      \scriptsize
      \begin{tabularx}{\linewidth}{@{}l L@{}}
        \toprule
        \cabecalho{Paradigma} & \cabecalho{Ideia central} \\
        \midrule
        Imperativo & sequência de comandos que alteram o estado (\emph{faça isto, depois aquilo}) \\
        Orientado a objetos & dados e comportamentos agrupados em objetos que trocam mensagens \\
        Funcional & funções puras, valores imutáveis, composição de funções \\
        Declarativo & descreve \textbf{o que} se quer, não \textbf{como} obter (SQL, interfaces em Compose e Flutter) \\
        \bottomrule
      \end{tabularx}
    \end{column}
    \begin{column}{0.44\textwidth}
      \begin{definicao}[Paradigma]
        Estilo de organizar o raciocínio e o código. Linguagens modernas são \textbf{multiparadigma}: Kotlin e Dart misturam imperativo, objetos e funcional.
      \end{definicao}
      \begin{dica}
        Use o paradigma que deixa cada parte mais clara: objetos para modelar o domínio, funções para transformar dados, declarativo para descrever a interface.
      \end{dica}
    \end{column}
  \end{columns}
  \begin{historia}
    Fortran (1957) era imperativo; Lisp (1958), funcional; Simula (1967) e Smalltalk (anos 1970) criaram a orientação a objetos. Kotlin (2011, JetBrains) e Dart (2011, Google) herdaram ideias de todos.
  \end{historia}
\end{frame}

\begin{frame}{Sistemas de tipos}
  \begin{columns}[T]
    \begin{column}{0.5\textwidth}
      \begin{definicao}[Tipo]
        Conjunto de valores possíveis e das operações permitidas sobre eles. Somar dois \cod{Int} faz sentido; somar um \cod{Int} a uma \cod{Bancada}, não.
      \end{definicao}
      \scriptsize
      \begin{tabularx}{\linewidth}{@{}l L@{}}
        \toprule
        \cabecalho{Classificação} & \cabecalho{Significado} \\
        \midrule
        Estática & tipos verificados na compilação (Kotlin, Dart, Java) \\
        Dinâmica & tipos verificados na execução (Python, JavaScript) \\
        Forte & conversões perigosas são proibidas ou explícitas \\
        Inferência & o compilador deduz o tipo (\cod{val x = 3}) \\
        \bottomrule
      \end{tabularx}
    \end{column}
    \begin{column}{0.48\textwidth}
      \begin{conceito}[Por que tipagem estática ajuda]
        Muitos erros aparecem \textbf{antes} de o app rodar; a IDE completa código e refatora com segurança; o tipo documenta a intenção.
      \end{conceito}
      \begin{pordentro}[Solidez (\emph{soundness})]
        Um sistema de tipos é sólido quando garante que, se o programa compila, uma variável nunca terá em execução um valor de outro tipo. O Dart 3 é sólido, inclusive quanto a nulos.
      \end{pordentro}
    \end{column}
  \end{columns}
\end{frame}

\begin{frame}{O ``erro de um bilhão de dólares'': o nulo}
  \begin{columns}[T]
    \begin{column}{0.5\textwidth}
      \begin{historia}
        Tony Hoare introduziu a referência nula na linguagem ALGOL W em 1965. Décadas depois, chamou-a de seu ``erro de um bilhão de dólares'', pelo custo das falhas que ela causou.
      \end{historia}
      \begin{definicao}[\emph{Null safety}]
        O sistema de tipos separa tipos que \textbf{podem} ser nulos (\cod{String?}) dos que \textbf{não podem} (\cod{String}). Usar um valor possivelmente nulo sem tratar vira erro de compilação, não falha em produção.
      \end{definicao}
    \end{column}
    \begin{column}{0.48\textwidth}
      \centering
      \begin{tikzpicture}[node distance=3mm,every node/.style={font=\scriptsize,align=center}]
        \node[caixa,minimum width=2.2cm] (a) {\cod{String?}\\ \tiny pode ser nulo};
        \node[caixapetroleo,right=8mm of a,minimum width=2.2cm] (b) {\cod{String}\\ \tiny nunca nulo};
        \draw[seta] (a) -- node[rotulo,above]{\tiny verificação} (b);
        \node[below=3mm of a,font=\tiny,text width=2.2cm,align=center] {\cod{if (x != null)}\\ \cod{x ?: padrão}};
      \end{tikzpicture}
      \begin{pordentro}[Análise de fluxo]
        Depois de um \cod{if (x != null)}, o compilador sabe que \cod{x} não é nulo naquele trecho e permite usá-lo diretamente (\emph{smart cast} no Kotlin, promoção de tipo no Dart).
      \end{pordentro}
      \begin{atencao}
        O operador \cod{!!} (Kotlin) ou \cod{!} (Dart) desliga a proteção: só use quando tiver certeza.
      \end{atencao}
    \end{column}
  \end{columns}
\end{frame}

\begin{frame}{Os pilares da orientação a objetos}
  \begin{columns}[T]
    \begin{column}{0.54\textwidth}
      \scriptsize
      \begin{tabularx}{\linewidth}{@{}l L@{}}
        \toprule
        \cabecalho{Pilar} & \cabecalho{Significado e exemplo} \\
        \midrule
        Encapsulamento & esconder o estado interno e expor operações: a \cod{Reserva} só muda pelo método \cod{cancelar()} \\
        Abstração & mostrar o essencial e ocultar detalhes: \cod{FonteLeituras} sem dizer se é BLE ou MQTT \\
        Herança & uma classe especializa outra: \cod{SensorTemperatura} herda de \cod{Sensor} \\
        Polimorfismo & o mesmo comando, comportamentos diferentes: \cod{fonte.leituras()} em qualquer implementação \\
        \bottomrule
      \end{tabularx}
    \end{column}
    \begin{column}{0.44\textwidth}
      \begin{definicao}[Objeto e classe]
        A \textbf{classe} é o molde (atributos e métodos); o \textbf{objeto} é uma instância concreta, com estado próprio.
      \end{definicao}
      \begin{dica}[Composição em vez de herança]
        Prefira montar objetos com outros objetos (``tem um'') a criar hierarquias profundas (``é um''): mudanças ficam localizadas e o código, mais flexível.
      \end{dica}
    \end{column}
  \end{columns}
  \begin{analogia}
    Uma tomada é uma abstração: qualquer aparelho com o plugue certo funciona (polimorfismo), e ninguém precisa conhecer a fiação dentro da parede (encapsulamento).
  \end{analogia}
\end{frame}

\begin{frame}{Programação funcional}
  \begin{columns}[T]
    \begin{column}{0.5\textwidth}
      \begin{definicao}[Função pura]
        Para a mesma entrada, devolve sempre a mesma saída e não causa efeitos colaterais (não altera variáveis externas, não grava em disco, não acessa a rede).
      \end{definicao}
      \begin{definicao}[Funções de primeira classe e de ordem superior]
        Funções são valores: podem ser guardadas em variáveis, passadas como argumento e devolvidas. Uma função que recebe ou devolve outra é de \textbf{ordem superior} (\cod{map}, \cod{filter}).
      \end{definicao}
    \end{column}
    \begin{column}{0.48\textwidth}
      \centering
      \begin{tikzpicture}[node distance=4mm,every node/.style={font=\scriptsize,align=center}]
        \node[caixa,minimum width=1.9cm] (l) {leituras};
        \node[caixapetroleo,right=of l,minimum width=1.6cm] (f) {\cod{filter}\\ \tiny > 30\,°C};
        \node[caixapetroleo,right=of f,minimum width=1.6cm] (m) {\cod{map}\\ \tiny para alerta};
        \node[caixaambar,below=4mm of m,minimum width=1.6cm] (r) {alertas};
        \draw[seta] (l) -- (f); \draw[seta] (f) -- (m); \draw[seta] (m) -- (r);
      \end{tikzpicture}
      \begin{conceito}[Por que importa no mobile]
        Funções puras são fáceis de testar e seguras entre \emph{threads}; transformar listas com \cod{map}/\cod{filter} é mais legível que laços; interfaces declarativas tratam a tela como uma função do estado (Módulo 05).
      \end{conceito}
    \end{column}
  \end{columns}
\end{frame}

\begin{frame}{Memória: pilha, \emph{heap} e coleta de lixo}
  \begin{columns}[c]
    \begin{column}{0.46\textwidth}
      \centering
      \begin{tikzpicture}[every node/.style={font=\tiny}]
        \node[font=\scriptsize\bfseries,text=azulpetroleo] at (0.9,2.7) {Pilha};
        \foreach \y/\t in {2.2/main(),1.7/carregar(),1.2/calcular()} \node[caixa,minimum width=1.8cm] at (0.9,\y) {\t};
        \node[font=\scriptsize\bfseries,text=azulpetroleo] at (3.7,2.7) {\emph{Heap}};
        \draw[cinzamedio,rounded corners] (2.6,0.3) rectangle (4.8,2.45);
        \node[caixapetroleo,minimum width=1.3cm] (o1) at (3.3,1.95) {Bancada};
        \node[caixapetroleo,minimum width=1.3cm] (o2) at (4.1,1.2) {Lista};
        \node[caixa,minimum width=1.3cm,fill=cinzaclaro] (o3) at (3.3,0.65) {sem ref.};
        \draw[seta] (1.8,2.2) -- (o1); \draw[seta] (1.8,1.7) -- (o2); \draw[seta] (o2) -- (o1);
        \node[font=\tiny,text=vermelhoerro] at (3.7,0.2) {coletado};
      \end{tikzpicture}
    \end{column}
    \begin{column}{0.52\textwidth}
      \begin{definicao}[Pilha e \emph{heap}]
        A \textbf{pilha} guarda as chamadas de função e suas variáveis locais, liberadas ao retornar. O \textbf{\emph{heap}} guarda os objetos, que vivem enquanto houver referências a eles.
      \end{definicao}
      \begin{definicao}[Coleta de lixo]
        O \emph{runtime} encontra objetos que não são mais alcançáveis a partir das referências vivas e libera sua memória automaticamente.
      \end{definicao}
      \begin{pordentro}[Gerações]
        A maioria dos objetos morre jovem. Por isso o ART e a VM do Dart separam objetos novos e antigos e coletam os novos com frequência e rapidez.
      \end{pordentro}
    \end{column}
  \end{columns}
\end{frame}

\begin{frame}{Imutabilidade e igualdade}
  \begin{columns}[T]
    \begin{column}{0.5\textwidth}
      \begin{definicao}[Imutável]
        Objeto que não muda depois de criado. Para ``alterar'', cria-se uma cópia com o campo novo (\cod{copy} no Kotlin, \cod{copyWith} no Dart).
      \end{definicao}
      \begin{conceito}[Vantagens]
        \begin{itemize}
          \item Ninguém altera o objeto ``pelas costas''.
          \item Seguro para compartilhar entre \emph{threads} e \emph{isolates}.
          \item Comparar o antes e o depois de um estado fica trivial (Módulo 07).
        \end{itemize}
      \end{conceito}
    \end{column}
    \begin{column}{0.48\textwidth}
      \begin{definicao}[Igualdade de valor e de identidade]
        \textbf{Identidade}: as duas variáveis apontam para o \textbf{mesmo} objeto. \textbf{Valor}: objetos diferentes com o mesmo conteúdo. \cod{data class} (Kotlin) e \emph{records} (Dart) comparam por valor.
      \end{definicao}
      \begin{erro}[Lista mutável compartilhada]
        Dois componentes recebem a mesma lista; um a altera e o outro mostra dados que não esperava. Exponha listas somente leitura.
      \end{erro}
    \end{column}
  \end{columns}
\end{frame}

% =============================================================================
\section{Fundamentos: linguagens de programação}

\begin{frame}{Paradigmas de programação}
  \begin{columns}[T]
    \begin{column}{0.54\textwidth}
      \scriptsize
      \begin{tabularx}{\linewidth}{@{}l L@{}}
        \toprule
        \cabecalho{Paradigma} & \cabecalho{Ideia central} \\
        \midrule
        Imperativo & sequência de comandos que alteram o estado (\emph{faça isto, depois aquilo}) \\
        Orientado a objetos & dados e comportamentos agrupados em objetos que trocam mensagens \\
        Funcional & funções puras, valores imutáveis, composição de funções \\
        Declarativo & descreve \textbf{o que} se quer, não \textbf{como} obter (SQL, interfaces em Compose e Flutter) \\
        \bottomrule
      \end{tabularx}
    \end{column}
    \begin{column}{0.44\textwidth}
      \begin{definicao}[Paradigma]
        Estilo de organizar o raciocínio e o código. Linguagens modernas são \textbf{multiparadigma}: Kotlin e Dart misturam imperativo, objetos e funcional.
      \end{definicao}
      \begin{dica}
        Use o paradigma que deixa cada parte mais clara: objetos para modelar o domínio, funções para transformar dados, declarativo para descrever a interface.
      \end{dica}
    \end{column}
  \end{columns}
  \begin{historia}
    Fortran (1957) era imperativo; Lisp (1958), funcional; Simula (1967) e Smalltalk (anos 1970) criaram a orientação a objetos. Kotlin (2011, JetBrains) e Dart (2011, Google) herdaram ideias de todos.
  \end{historia}
\end{frame}

\begin{frame}{Sistemas de tipos}
  \begin{columns}[T]
    \begin{column}{0.5\textwidth}
      \begin{definicao}[Tipo]
        Conjunto de valores possíveis e das operações permitidas sobre eles. Somar dois \cod{Int} faz sentido; somar um \cod{Int} a uma \cod{Bancada}, não.
      \end{definicao}
      \scriptsize
      \begin{tabularx}{\linewidth}{@{}l L@{}}
        \toprule
        \cabecalho{Classificação} & \cabecalho{Significado} \\
        \midrule
        Estática & tipos verificados na compilação (Kotlin, Dart, Java) \\
        Dinâmica & tipos verificados na execução (Python, JavaScript) \\
        Forte & conversões perigosas são proibidas ou explícitas \\
        Inferência & o compilador deduz o tipo (\cod{val x = 3}) \\
        \bottomrule
      \end{tabularx}
    \end{column}
    \begin{column}{0.48\textwidth}
      \begin{conceito}[Por que tipagem estática ajuda]
        Muitos erros aparecem \textbf{antes} de o app rodar; a IDE completa código e refatora com segurança; o tipo documenta a intenção.
      \end{conceito}
      \begin{pordentro}[Solidez (\emph{soundness})]
        Um sistema de tipos é sólido quando garante que, se o programa compila, uma variável nunca terá em execução um valor de outro tipo. O Dart 3 é sólido, inclusive quanto a nulos.
      \end{pordentro}
    \end{column}
  \end{columns}
\end{frame}

\begin{frame}{O ``erro de um bilhão de dólares'': o nulo}
  \begin{columns}[T]
    \begin{column}{0.5\textwidth}
      \begin{historia}
        Tony Hoare introduziu a referência nula na linguagem ALGOL W em 1965. Décadas depois, chamou-a de seu ``erro de um bilhão de dólares'', pelo custo das falhas que ela causou.
      \end{historia}
      \begin{definicao}[\emph{Null safety}]
        O sistema de tipos separa tipos que \textbf{podem} ser nulos (\cod{String?}) dos que \textbf{não podem} (\cod{String}). Usar um valor possivelmente nulo sem tratar vira erro de compilação, não falha em produção.
      \end{definicao}
    \end{column}
    \begin{column}{0.48\textwidth}
      \centering
      \begin{tikzpicture}[node distance=3mm,every node/.style={font=\scriptsize,align=center}]
        \node[caixa,minimum width=2.2cm] (a) {\cod{String?}\\ \tiny pode ser nulo};
        \node[caixapetroleo,right=8mm of a,minimum width=2.2cm] (b) {\cod{String}\\ \tiny nunca nulo};
        \draw[seta] (a) -- node[rotulo,above]{\tiny verificação} (b);
        \node[below=3mm of a,font=\tiny,text width=2.2cm,align=center] {\cod{if (x != null)}\\ \cod{x ?: padrão}};
      \end{tikzpicture}
      \begin{pordentro}[Análise de fluxo]
        Depois de um \cod{if (x != null)}, o compilador sabe que \cod{x} não é nulo naquele trecho e permite usá-lo diretamente (\emph{smart cast} no Kotlin, promoção de tipo no Dart).
      \end{pordentro}
      \begin{atencao}
        O operador \cod{!!} (Kotlin) ou \cod{!} (Dart) desliga a proteção: só use quando tiver certeza.
      \end{atencao}
    \end{column}
  \end{columns}
\end{frame}

\begin{frame}{Os pilares da orientação a objetos}
  \begin{columns}[T]
    \begin{column}{0.54\textwidth}
      \scriptsize
      \begin{tabularx}{\linewidth}{@{}l L@{}}
        \toprule
        \cabecalho{Pilar} & \cabecalho{Significado e exemplo} \\
        \midrule
        Encapsulamento & esconder o estado interno e expor operações: a \cod{Reserva} só muda pelo método \cod{cancelar()} \\
        Abstração & mostrar o essencial e ocultar detalhes: \cod{FonteLeituras} sem dizer se é BLE ou MQTT \\
        Herança & uma classe especializa outra: \cod{SensorTemperatura} herda de \cod{Sensor} \\
        Polimorfismo & o mesmo comando, comportamentos diferentes: \cod{fonte.leituras()} em qualquer implementação \\
        \bottomrule
      \end{tabularx}
    \end{column}
    \begin{column}{0.44\textwidth}
      \begin{definicao}[Objeto e classe]
        A \textbf{classe} é o molde (atributos e métodos); o \textbf{objeto} é uma instância concreta, com estado próprio.
      \end{definicao}
      \begin{dica}[Composição em vez de herança]
        Prefira montar objetos com outros objetos (``tem um'') a criar hierarquias profundas (``é um''): mudanças ficam localizadas e o código, mais flexível.
      \end{dica}
    \end{column}
  \end{columns}
  \begin{analogia}
    Uma tomada é uma abstração: qualquer aparelho com o plugue certo funciona (polimorfismo), e ninguém precisa conhecer a fiação dentro da parede (encapsulamento).
  \end{analogia}
\end{frame}

\begin{frame}{Programação funcional}
  \begin{columns}[T]
    \begin{column}{0.5\textwidth}
      \begin{definicao}[Função pura]
        Para a mesma entrada, devolve sempre a mesma saída e não causa efeitos colaterais (não altera variáveis externas, não grava em disco, não acessa a rede).
      \end{definicao}
      \begin{definicao}[Funções de primeira classe e de ordem superior]
        Funções são valores: podem ser guardadas em variáveis, passadas como argumento e devolvidas. Uma função que recebe ou devolve outra é de \textbf{ordem superior} (\cod{map}, \cod{filter}).
      \end{definicao}
    \end{column}
    \begin{column}{0.48\textwidth}
      \centering
      \begin{tikzpicture}[node distance=4mm,every node/.style={font=\scriptsize,align=center}]
        \node[caixa,minimum width=1.9cm] (l) {leituras};
        \node[caixapetroleo,right=of l,minimum width=1.6cm] (f) {\cod{filter}\\ \tiny > 30\,°C};
        \node[caixapetroleo,right=of f,minimum width=1.6cm] (m) {\cod{map}\\ \tiny para alerta};
        \node[caixaambar,below=4mm of m,minimum width=1.6cm] (r) {alertas};
        \draw[seta] (l) -- (f); \draw[seta] (f) -- (m); \draw[seta] (m) -- (r);
      \end{tikzpicture}
      \begin{conceito}[Por que importa no mobile]
        Funções puras são fáceis de testar e seguras entre \emph{threads}; transformar listas com \cod{map}/\cod{filter} é mais legível que laços; interfaces declarativas tratam a tela como uma função do estado (Módulo 05).
      \end{conceito}
    \end{column}
  \end{columns}
\end{frame}

\begin{frame}{Memória: pilha, \emph{heap} e coleta de lixo}
  \begin{columns}[c]
    \begin{column}{0.46\textwidth}
      \centering
      \begin{tikzpicture}[every node/.style={font=\tiny}]
        \node[font=\scriptsize\bfseries,text=azulpetroleo] at (0.9,2.7) {Pilha};
        \foreach \y/\t in {2.2/main(),1.7/carregar(),1.2/calcular()} \node[caixa,minimum width=1.8cm] at (0.9,\y) {\t};
        \node[font=\scriptsize\bfseries,text=azulpetroleo] at (3.7,2.7) {\emph{Heap}};
        \draw[cinzamedio,rounded corners] (2.6,0.3) rectangle (4.8,2.45);
        \node[caixapetroleo,minimum width=1.3cm] (o1) at (3.3,1.95) {Bancada};
        \node[caixapetroleo,minimum width=1.3cm] (o2) at (4.1,1.2) {Lista};
        \node[caixa,minimum width=1.3cm,fill=cinzaclaro] (o3) at (3.3,0.65) {sem ref.};
        \draw[seta] (1.8,2.2) -- (o1); \draw[seta] (1.8,1.7) -- (o2); \draw[seta] (o2) -- (o1);
        \node[font=\tiny,text=vermelhoerro] at (3.7,0.2) {coletado};
      \end{tikzpicture}
    \end{column}
    \begin{column}{0.52\textwidth}
      \begin{definicao}[Pilha e \emph{heap}]
        A \textbf{pilha} guarda as chamadas de função e suas variáveis locais, liberadas ao retornar. O \textbf{\emph{heap}} guarda os objetos, que vivem enquanto houver referências a eles.
      \end{definicao}
      \begin{definicao}[Coleta de lixo]
        O \emph{runtime} encontra objetos que não são mais alcançáveis a partir das referências vivas e libera sua memória automaticamente.
      \end{definicao}
      \begin{pordentro}[Gerações]
        A maioria dos objetos morre jovem. Por isso o ART e a VM do Dart separam objetos novos e antigos e coletam os novos com frequência e rapidez.
      \end{pordentro}
    \end{column}
  \end{columns}
\end{frame}

\begin{frame}{Imutabilidade e igualdade}
  \begin{columns}[T]
    \begin{column}{0.5\textwidth}
      \begin{definicao}[Imutável]
        Objeto que não muda depois de criado. Para ``alterar'', cria-se uma cópia com o campo novo (\cod{copy} no Kotlin, \cod{copyWith} no Dart).
      \end{definicao}
      \begin{conceito}[Vantagens]
        \begin{itemize}
          \item Ninguém altera o objeto ``pelas costas''.
          \item Seguro para compartilhar entre \emph{threads} e \emph{isolates}.
          \item Comparar o antes e o depois de um estado fica trivial (Módulo 07).
        \end{itemize}
      \end{conceito}
    \end{column}
    \begin{column}{0.48\textwidth}
      \begin{definicao}[Igualdade de valor e de identidade]
        \textbf{Identidade}: as duas variáveis apontam para o \textbf{mesmo} objeto. \textbf{Valor}: objetos diferentes com o mesmo conteúdo. \cod{data class} (Kotlin) e \emph{records} (Dart) comparam por valor.
      \end{definicao}
      \begin{erro}[Lista mutável compartilhada]
        Dois componentes recebem a mesma lista; um a altera e o outro mostra dados que não esperava. Exponha listas somente leitura.
      \end{erro}
    \end{column}
  \end{columns}
\end{frame}

% =============================================================================
\section{Fundamentos: linguagens de programação}

\begin{frame}{Paradigmas de programação}
  \begin{columns}[T]
    \begin{column}{0.54\textwidth}
      \scriptsize
      \begin{tabularx}{\linewidth}{@{}l L@{}}
        \toprule
        \cabecalho{Paradigma} & \cabecalho{Ideia central} \\
        \midrule
        Imperativo & sequência de comandos que alteram o estado (\emph{faça isto, depois aquilo}) \\
        Orientado a objetos & dados e comportamentos agrupados em objetos que trocam mensagens \\
        Funcional & funções puras, valores imutáveis, composição de funções \\
        Declarativo & descreve \textbf{o que} se quer, não \textbf{como} obter (SQL, interfaces em Compose e Flutter) \\
        \bottomrule
      \end{tabularx}
    \end{column}
    \begin{column}{0.44\textwidth}
      \begin{definicao}[Paradigma]
        Estilo de organizar o raciocínio e o código. Linguagens modernas são \textbf{multiparadigma}: Kotlin e Dart misturam imperativo, objetos e funcional.
      \end{definicao}
      \begin{dica}
        Use o paradigma que deixa cada parte mais clara: objetos para modelar o domínio, funções para transformar dados, declarativo para descrever a interface.
      \end{dica}
    \end{column}
  \end{columns}
  \begin{historia}
    Fortran (1957) era imperativo; Lisp (1958), funcional; Simula (1967) e Smalltalk (anos 1970) criaram a orientação a objetos. Kotlin (2011, JetBrains) e Dart (2011, Google) herdaram ideias de todos.
  \end{historia}
\end{frame}

\begin{frame}{Sistemas de tipos}
  \begin{columns}[T]
    \begin{column}{0.5\textwidth}
      \begin{definicao}[Tipo]
        Conjunto de valores possíveis e das operações permitidas sobre eles. Somar dois \cod{Int} faz sentido; somar um \cod{Int} a uma \cod{Bancada}, não.
      \end{definicao}
      \scriptsize
      \begin{tabularx}{\linewidth}{@{}l L@{}}
        \toprule
        \cabecalho{Classificação} & \cabecalho{Significado} \\
        \midrule
        Estática & tipos verificados na compilação (Kotlin, Dart, Java) \\
        Dinâmica & tipos verificados na execução (Python, JavaScript) \\
        Forte & conversões perigosas são proibidas ou explícitas \\
        Inferência & o compilador deduz o tipo (\cod{val x = 3}) \\
        \bottomrule
      \end{tabularx}
    \end{column}
    \begin{column}{0.48\textwidth}
      \begin{conceito}[Por que tipagem estática ajuda]
        Muitos erros aparecem \textbf{antes} de o app rodar; a IDE completa código e refatora com segurança; o tipo documenta a intenção.
      \end{conceito}
      \begin{pordentro}[Solidez (\emph{soundness})]
        Um sistema de tipos é sólido quando garante que, se o programa compila, uma variável nunca terá em execução um valor de outro tipo. O Dart 3 é sólido, inclusive quanto a nulos.
      \end{pordentro}
    \end{column}
  \end{columns}
\end{frame}

\begin{frame}{O ``erro de um bilhão de dólares'': o nulo}
  \begin{columns}[T]
    \begin{column}{0.5\textwidth}
      \begin{historia}
        Tony Hoare introduziu a referência nula na linguagem ALGOL W em 1965. Décadas depois, chamou-a de seu ``erro de um bilhão de dólares'', pelo custo das falhas que ela causou.
      \end{historia}
      \begin{definicao}[\emph{Null safety}]
        O sistema de tipos separa tipos que \textbf{podem} ser nulos (\cod{String?}) dos que \textbf{não podem} (\cod{String}). Usar um valor possivelmente nulo sem tratar vira erro de compilação, não falha em produção.
      \end{definicao}
    \end{column}
    \begin{column}{0.48\textwidth}
      \centering
      \begin{tikzpicture}[node distance=3mm,every node/.style={font=\scriptsize,align=center}]
        \node[caixa,minimum width=2.2cm] (a) {\cod{String?}\\ \tiny pode ser nulo};
        \node[caixapetroleo,right=8mm of a,minimum width=2.2cm] (b) {\cod{String}\\ \tiny nunca nulo};
        \draw[seta] (a) -- node[rotulo,above]{\tiny verificação} (b);
        \node[below=3mm of a,font=\tiny,text width=2.2cm,align=center] {\cod{if (x != null)}\\ \cod{x ?: padrão}};
      \end{tikzpicture}
      \begin{pordentro}[Análise de fluxo]
        Depois de um \cod{if (x != null)}, o compilador sabe que \cod{x} não é nulo naquele trecho e permite usá-lo diretamente (\emph{smart cast} no Kotlin, promoção de tipo no Dart).
      \end{pordentro}
      \begin{atencao}
        O operador \cod{!!} (Kotlin) ou \cod{!} (Dart) desliga a proteção: só use quando tiver certeza.
      \end{atencao}
    \end{column}
  \end{columns}
\end{frame}

\begin{frame}{Os pilares da orientação a objetos}
  \begin{columns}[T]
    \begin{column}{0.54\textwidth}
      \scriptsize
      \begin{tabularx}{\linewidth}{@{}l L@{}}
        \toprule
        \cabecalho{Pilar} & \cabecalho{Significado e exemplo} \\
        \midrule
        Encapsulamento & esconder o estado interno e expor operações: a \cod{Reserva} só muda pelo método \cod{cancelar()} \\
        Abstração & mostrar o essencial e ocultar detalhes: \cod{FonteLeituras} sem dizer se é BLE ou MQTT \\
        Herança & uma classe especializa outra: \cod{SensorTemperatura} herda de \cod{Sensor} \\
        Polimorfismo & o mesmo comando, comportamentos diferentes: \cod{fonte.leituras()} em qualquer implementação \\
        \bottomrule
      \end{tabularx}
    \end{column}
    \begin{column}{0.44\textwidth}
      \begin{definicao}[Objeto e classe]
        A \textbf{classe} é o molde (atributos e métodos); o \textbf{objeto} é uma instância concreta, com estado próprio.
      \end{definicao}
      \begin{dica}[Composição em vez de herança]
        Prefira montar objetos com outros objetos (``tem um'') a criar hierarquias profundas (``é um''): mudanças ficam localizadas e o código, mais flexível.
      \end{dica}
    \end{column}
  \end{columns}
  \begin{analogia}
    Uma tomada é uma abstração: qualquer aparelho com o plugue certo funciona (polimorfismo), e ninguém precisa conhecer a fiação dentro da parede (encapsulamento).
  \end{analogia}
\end{frame}

\begin{frame}{Programação funcional}
  \begin{columns}[T]
    \begin{column}{0.5\textwidth}
      \begin{definicao}[Função pura]
        Para a mesma entrada, devolve sempre a mesma saída e não causa efeitos colaterais (não altera variáveis externas, não grava em disco, não acessa a rede).
      \end{definicao}
      \begin{definicao}[Funções de primeira classe e de ordem superior]
        Funções são valores: podem ser guardadas em variáveis, passadas como argumento e devolvidas. Uma função que recebe ou devolve outra é de \textbf{ordem superior} (\cod{map}, \cod{filter}).
      \end{definicao}
    \end{column}
    \begin{column}{0.48\textwidth}
      \centering
      \begin{tikzpicture}[node distance=4mm,every node/.style={font=\scriptsize,align=center}]
        \node[caixa,minimum width=1.9cm] (l) {leituras};
        \node[caixapetroleo,right=of l,minimum width=1.6cm] (f) {\cod{filter}\\ \tiny > 30\,°C};
        \node[caixapetroleo,right=of f,minimum width=1.6cm] (m) {\cod{map}\\ \tiny para alerta};
        \node[caixaambar,below=4mm of m,minimum width=1.6cm] (r) {alertas};
        \draw[seta] (l) -- (f); \draw[seta] (f) -- (m); \draw[seta] (m) -- (r);
      \end{tikzpicture}
      \begin{conceito}[Por que importa no mobile]
        Funções puras são fáceis de testar e seguras entre \emph{threads}; transformar listas com \cod{map}/\cod{filter} é mais legível que laços; interfaces declarativas tratam a tela como uma função do estado (Módulo 05).
      \end{conceito}
    \end{column}
  \end{columns}
\end{frame}

\begin{frame}{Memória: pilha, \emph{heap} e coleta de lixo}
  \begin{columns}[c]
    \begin{column}{0.46\textwidth}
      \centering
      \begin{tikzpicture}[every node/.style={font=\tiny}]
        \node[font=\scriptsize\bfseries,text=azulpetroleo] at (0.9,2.7) {Pilha};
        \foreach \y/\t in {2.2/main(),1.7/carregar(),1.2/calcular()} \node[caixa,minimum width=1.8cm] at (0.9,\y) {\t};
        \node[font=\scriptsize\bfseries,text=azulpetroleo] at (3.7,2.7) {\emph{Heap}};
        \draw[cinzamedio,rounded corners] (2.6,0.3) rectangle (4.8,2.45);
        \node[caixapetroleo,minimum width=1.3cm] (o1) at (3.3,1.95) {Bancada};
        \node[caixapetroleo,minimum width=1.3cm] (o2) at (4.1,1.2) {Lista};
        \node[caixa,minimum width=1.3cm,fill=cinzaclaro] (o3) at (3.3,0.65) {sem ref.};
        \draw[seta] (1.8,2.2) -- (o1); \draw[seta] (1.8,1.7) -- (o2); \draw[seta] (o2) -- (o1);
        \node[font=\tiny,text=vermelhoerro] at (3.7,0.2) {coletado};
      \end{tikzpicture}
    \end{column}
    \begin{column}{0.52\textwidth}
      \begin{definicao}[Pilha e \emph{heap}]
        A \textbf{pilha} guarda as chamadas de função e suas variáveis locais, liberadas ao retornar. O \textbf{\emph{heap}} guarda os objetos, que vivem enquanto houver referências a eles.
      \end{definicao}
      \begin{definicao}[Coleta de lixo]
        O \emph{runtime} encontra objetos que não são mais alcançáveis a partir das referências vivas e libera sua memória automaticamente.
      \end{definicao}
      \begin{pordentro}[Gerações]
        A maioria dos objetos morre jovem. Por isso o ART e a VM do Dart separam objetos novos e antigos e coletam os novos com frequência e rapidez.
      \end{pordentro}
    \end{column}
  \end{columns}
\end{frame}

\begin{frame}{Imutabilidade e igualdade}
  \begin{columns}[T]
    \begin{column}{0.5\textwidth}
      \begin{definicao}[Imutável]
        Objeto que não muda depois de criado. Para ``alterar'', cria-se uma cópia com o campo novo (\cod{copy} no Kotlin, \cod{copyWith} no Dart).
      \end{definicao}
      \begin{conceito}[Vantagens]
        \begin{itemize}
          \item Ninguém altera o objeto ``pelas costas''.
          \item Seguro para compartilhar entre \emph{threads} e \emph{isolates}.
          \item Comparar o antes e o depois de um estado fica trivial (Módulo 07).
        \end{itemize}
      \end{conceito}
    \end{column}
    \begin{column}{0.48\textwidth}
      \begin{definicao}[Igualdade de valor e de identidade]
        \textbf{Identidade}: as duas variáveis apontam para o \textbf{mesmo} objeto. \textbf{Valor}: objetos diferentes com o mesmo conteúdo. \cod{data class} (Kotlin) e \emph{records} (Dart) comparam por valor.
      \end{definicao}
      \begin{erro}[Lista mutável compartilhada]
        Dois componentes recebem a mesma lista; um a altera e o outro mostra dados que não esperava. Exponha listas somente leitura.
      \end{erro}
    \end{column}
  \end{columns}
\end{frame}


% =============================================================================
\section{Primeiros passos}

\begin{frame}[fragile]{Variáveis e tipos}
  \begin{columns}[T]
    \begin{column}{0.54\textwidth}
\begin{codigo}[title=variaveis.dart]{Dart}
void main() {
  var bancada = 3;               // tipo inferido: int
  double temperatura = 24.5;
  final nome = 'Osciloscópio';   // atribui uma vez
  const limite = 30.0;           // constante de compilação
  bool ligado = true;
  late String relatorio;         // inicializada depois

  temperatura += 0.3;
  relatorio = 'Bancada $bancada: $temperatura °C';
  print(relatorio);
  print('Nome tem ${nome.length} letras');

  final texto = '''
Equipamento: $nome
Estado: ${ligado ? 'ON' : 'OFF'}''';
}
\end{codigo}
    \end{column}
    \begin{column}{0.44\textwidth}
      \scriptsize
      \begin{tabularx}{\linewidth}{@{}l L@{}}
        \toprule
        \cabecalho{Palavra} & \cabecalho{Significado} \\
        \midrule
        \cod{var} & mutável, tipo inferido \\
        \cod{final} & atribuída uma única vez, em execução \\
        \cod{const} & valor fixo em compilação \\
        \cod{late} & inicializada depois, antes do uso \\
        \bottomrule
      \end{tabularx}
      \begin{conceito}[Tipos básicos]
        \cod{int}, \cod{double}, \cod{num}, \cod{String}, \cod{bool}, além de \cod{List}, \cod{Set}, \cod{Map} e \cod{Object}. Tudo em Dart é objeto.
      \end{conceito}
      \begin{dica}
        Prefira \cod{final}; use \cod{var} só para o que muda.
      \end{dica}
    \end{column}
  \end{columns}
\end{frame}

\begin{frame}[fragile]{Null safety}
  \begin{columns}[T]
    \begin{column}{0.54\textwidth}
\begin{codigo}[title=null\_safety.dart]{Dart}
String nome = 'LDR';
// nome = null;            // ERRO de compilação

String? apelido;           // pode ser nulo
// apelido.length;         // ERRO: pode ser nulo

int? tam = apelido?.length;          // chamada segura
int tam2 = apelido?.length ?? 0;     // valor padrão
apelido ??= 'sem apelido';           // atribui se nulo

if (apelido != null) {
  print(apelido.length);   // promovido a String
}

final tam3 = apelido!.length;  // "garanto que não é
                               // nulo": evite!
\end{codigo}
    \end{column}
    \begin{column}{0.44\textwidth}
      \begin{conceito}[A ideia]
        \cod{String} nunca é nulo; \cod{String?} pode ser. O compilador obriga a tratar o caso nulo antes de usar o valor.
      \end{conceito}
      \scriptsize
      \begin{tabular}{@{}l l@{}}
        \toprule
        \cabecalho{Operador} & \cabecalho{Uso} \\
        \midrule
        \cod{?.} & acessa se não for nulo \\
        \cod{??} & valor padrão \\
        \cod{??=} & atribui se for nulo \\
        \cod{!} & afirma não nulo (lança erro se for) \\
        \bottomrule
      \end{tabular}
      \begin{erro}
        Espalhar \cod{!} pelo código traz de volta os erros de nulo que o Dart tenta evitar.
      \end{erro}
    \end{column}
  \end{columns}
\end{frame}

\begin{frame}[fragile]{Controle de fluxo}
  \begin{columns}[T]
    \begin{column}{0.5\textwidth}
\begin{codigo}[title=Condições]{Dart}
final status = temp > 80 ? 'ALERTA' : 'OK';

// switch como expressão (Dart 3)
String classificar(double t) => switch (t) {
  < 0 => 'abaixo da faixa',
  >= 0 && <= 35 => 'normal',
  _ => 'acima da faixa',
};

String descrever(Object x) => switch (x) {
  int n => 'inteiro $n',
  String s => 'texto de ${s.length} letras',
  _ => 'desconhecido',
};
\end{codigo}
    \end{column}
    \begin{column}{0.48\textwidth}
\begin{codigo}[title=Laços]{Dart}
for (var i = 0; i < 5; i++) print(i);

final sensores = ['DHT22', 'BMP280', 'LDR'];
for (final s in sensores) print(s);
sensores.forEach(print);

var tentativas = 0;
while (tentativas < 3) {
  tentativas++;
}
\end{codigo}
      \begin{dica}
        O \cod{switch} de expressão do Dart 3 aceita \emph{patterns} (faixas, tipos, \emph{records}) e o compilador avisa quando faltam casos.
      \end{dica}
    \end{column}
  \end{columns}
\end{frame}

% =============================================================================
\section{Funções}

\begin{frame}[fragile]{Funções e parâmetros}
  \begin{columns}[T]
    \begin{column}{0.52\textwidth}
\begin{codigo}[title=funcoes.dart]{Dart}
double media(List<double> valores) {
  return valores.reduce((a, b) => a + b) / valores.length;
}

// corpo de expressão única (arrow)
double cParaF(double c) => c * 9 / 5 + 32;

// parâmetros nomeados (com required e padrão)
void conectar({
  required String host,
  int porta = 1883,
  bool tls = false,
}) { /* ... */ }

// parâmetros posicionais opcionais
String saudacao(String nome, [String? titulo]) =>
    'Olá, ${titulo ?? ''} $nome';

conectar(host: 'broker.local', tls: true);
\end{codigo}
    \end{column}
    \begin{column}{0.46\textwidth}
      \begin{conceito}[Parâmetros nomeados]
        Ficam entre \cod{\{ \}}, são opcionais por padrão e podem ser \cod{required}. É assim que todos os widgets do Flutter recebem seus parâmetros: \cod{Text('Oi', style: \ldots)}.
      \end{conceito}
      \begin{dica}
        A seta \cod{=>} substitui \cod{\{ return \ldots; \}} quando o corpo é uma única expressão.
      \end{dica}
    \end{column}
  \end{columns}
\end{frame}

\begin{frame}[fragile]{Funções como valores e extensões}
  \begin{columns}[T]
    \begin{column}{0.52\textwidth}
\begin{codigo}[title=Funções anônimas]{Dart}
// funções são objetos
int Function(int) dobro = (x) => x * 2;
final ehPar = (int n) => n % 2 == 0;

// recebe uma função como parâmetro
Duration medir(int n, void Function() acao) {
  final sw = Stopwatch()..start();
  for (var i = 0; i < n; i++) acao();
  return sw.elapsed;
}

final tempo = medir(1000, () => calcularFFT());
\end{codigo}
    \end{column}
    \begin{column}{0.46\textwidth}
\begin{codigo}[title=Extensões]{Dart}
extension Conversoes on double {
  String get emCelsius =>
      '${toStringAsFixed(1)} °C';
}
extension Adc on int {
  double paraVolts({double vref = 3.3,
                    int bits = 12}) =>
      this * vref / ((1 << bits) - 1);
}

print(24.567.emCelsius);   // 24.6 °C
print(2048.paraVolts());   // 1.65
\end{codigo}
      \begin{atencao}[Isto é Flutter!]
        \cod{onPressed: () \{ \ldots \}} é uma função anônima passada como parâmetro nomeado.
      \end{atencao}
    \end{column}
  \end{columns}
\end{frame}

% =============================================================================
\section{Coleções}

\begin{frame}[fragile]{List, Set e Map}
  \begin{columns}[T]
    \begin{column}{0.52\textwidth}
\begin{codigo}[title=Coleções]{Dart}
final nomes = <String>['DHT22', 'BMP280', 'LDR'];
final ids = {1, 2, 2, 3};               // Set: {1, 2, 3}
final pinos = {'led': 2, 'rele': 4};    // Map

nomes.add('DS18B20');
pinos['buzzer'] = 5;
final pinoLed = pinos['led'] ?? -1;     // pode ser nulo

const fixas = ['A', 'B'];   // lista imutável
// fixas.add('C');           // erro em execução

// collection if/for e spread
final menu = [
  'Painel',
  if (ehAdmin) 'Usuários',
  for (final b in bancadas) 'Bancada ${b.id}',
  ...extras,
];
\end{codigo}
    \end{column}
    \begin{column}{0.46\textwidth}
      \footnotesize
      \begin{tabular}{@{}l l@{}}
        \toprule
        \cabecalho{Tipo} & \cabecalho{Característica} \\
        \midrule
        \cod{List} & ordenada, aceita repetidos \\
        \cod{Set} & sem repetidos \\
        \cod{Map} & pares chave $\rightarrow$ valor \\
        \bottomrule
      \end{tabular}
      \begin{conceito}[\emph{Collection if/for}]
        \cod{if}, \cod{for} e \cod{...} (\emph{spread}) dentro de literais de lista. Muito usados para montar listas de widgets condicionalmente.
      \end{conceito}
    \end{column}
  \end{columns}
\end{frame}

\begin{frame}[fragile]{Operações funcionais}
  \begin{columns}[T]
    \begin{column}{0.56\textwidth}
\begin{codigo}[title=Pipeline de dados]{Dart}
final leituras = [
  Leitura('t1', 24.1), Leitura('t2', 81.3),
  Leitura('t1', 24.9), Leitura('t3', 19.0),
];

final alertas = leituras
    .where((l) => l.valor > 75)           // filtra
    .map((l) => '${l.sensor}: ${l.valor}') // transforma
    .toList();

final total = leituras.fold<double>(0, (s, l) => s + l.valor);
final temAlerta = leituras.any((l) => l.valor > 80);
final t3 = leituras.where((l) => l.sensor == 't3').firstOrNull;
final ordenadas = [...leituras]   // .. = cascata
  ..sort((a, b) => b.valor.compareTo(a.valor));
\end{codigo}
    \end{column}
    \begin{column}{0.42\textwidth}
      \scriptsize
      \begin{tabular}{@{}l l@{}}
        \toprule
        \cabecalho{Método} & \cabecalho{Resultado} \\
        \midrule
        \cod{where} & mantém os que passam \\
        \cod{map} & transforma cada item \\
        \cod{fold}, \cod{reduce} & agrega \\
        \cod{any}, \cod{every} & testes lógicos \\
        \cod{firstWhere} & primeiro que satisfaz \\
        \cod{sort} & ordena a própria lista \\
        \cod{toList}, \cod{toSet} & materializa \\
        \bottomrule
      \end{tabular}
      \begin{dica}
        \cod{where} e \cod{map} devolvem \cod{Iterable} preguiçoso: chame \cod{toList()} quando precisar de uma lista.
      \end{dica}
    \end{column}
  \end{columns}
\end{frame}

% =============================================================================
\section{Orientação a objetos}

\begin{frame}[fragile]{Classes e construtores}
  \begin{columns}[T]
    \begin{column}{0.54\textwidth}
\begin{codigo}[title=equipamento.dart]{Dart}
class Equipamento {
  final String patrimonio;
  String descricao;
  int _horasDeUso = 0;          // _ = privado ao arquivo

  Equipamento(this.patrimonio, this.descricao)
      : assert(patrimonio.isNotEmpty);

  // construtor nomeado
  Equipamento.semDescricao(String p) : this(p, '---');

  bool get precisaManutencao => _horasDeUso > 500;

  void registrarUso(int horas) => _horasDeUso += horas;
}

final osc = Equipamento('IFMA-0042', 'Osciloscópio');
osc.registrarUso(12);
\end{codigo}
    \end{column}
    \begin{column}{0.44\textwidth}
      \begin{itemize}
        \item \cod{this.campo} no construtor atribui direto.
        \item Lista de inicialização após \cod{:}.
        \item Construtores nomeados: \cod{Classe.nome()}.
        \item \cod{get}/\cod{set} para propriedades calculadas.
        \item Prefixo \cod{\_} = privado à biblioteca (arquivo).
        \item Não precisa de \cod{new}.
      \end{itemize}
      \begin{dica}[Construtores \cod{const} e \cod{factory}]
        \cod{const Classe(\ldots)} cria objetos imutáveis compartilhados (como os widgets). \cod{factory} pode devolver um objeto já existente ou de outro subtipo.
      \end{dica}
    \end{column}
  \end{columns}
\end{frame}

\begin{frame}[fragile]{Herança, interfaces, mixins}
  \begin{columns}[T]
    \begin{column}{0.54\textwidth}
\begin{codigo}[title=Hierarquia]{Dart}
abstract class Medidor {
  String get unidade;
  double ler();
  String formatado() => '${ler()} $unidade';
}

mixin Calibravel {
  double offset = 0;
  void calibrar(double referencia, double lido) =>
      offset = referencia - lido;
}

class Sensor { final int pino; Sensor(this.pino); }

class Termistor extends Sensor
    with Calibravel implements Medidor {
  Termistor(super.pino);
  @override String get unidade => '°C';
  @override double ler() => 25.0 + offset;
}
\end{codigo}
    \end{column}
    \begin{column}{0.44\textwidth}
      \begin{conceito}[Regras]
        \begin{itemize}
          \item \cod{extends}: herança simples.
          \item \cod{implements}: toda classe é também uma interface.
          \item \cod{with}: \emph{mixin} reaproveita código em várias classes.
          \item \cod{abstract}: não pode ser instanciada.
        \end{itemize}
      \end{conceito}
      \begin{dica}
        \cod{@override} documenta a sobrescrita e o analisador avisa se o método não existir na superclasse.
      \end{dica}
    \end{column}
  \end{columns}
\end{frame}

\begin{frame}[fragile]{Enum, records e patterns}
  \begin{columns}[T]
    \begin{column}{0.5\textwidth}
\begin{codigo}[title=Enum aprimorado]{Dart}
enum StatusBancada {
  livre('verde'),
  reservada('âmbar'),
  manutencao('vermelho');

  const StatusBancada(this.cor);
  final String cor;
  bool get disponivel => this == livre;
}

print(StatusBancada.values);
\end{codigo}
    \end{column}
    \begin{column}{0.48\textwidth}
\begin{codigo}[title=Records e desestruturação]{Dart}
// record: vários valores sem criar classe
(double, double) faixa() => (10.0, 40.0);
({double min, double max}) faixa2() =>
    (min: 10, max: 40);

final (min, max) = faixa();
final (:min2, :max2) = (min2: 1, max2: 2);

// pattern em JSON
if (json case {'id': int id, 'nome': String n}) {
  print('Bancada $id: $n');
}
\end{codigo}
    \end{column}
  \end{columns}
  \begin{conceito}[\emph{Records} e \emph{patterns} (Dart 3)]
    \emph{Records} agrupam valores com tipo, sem criar uma classe. \emph{Patterns} desestruturam e validam dados, muito úteis para ler JSON com segurança.
  \end{conceito}
\end{frame}

\begin{frame}[fragile]{Sealed classes: modelando estados}
  \begin{columns}[T]
    \begin{column}{0.54\textwidth}
\begin{codigo}[title=resultado.dart]{Dart}
sealed class Resultado {}
class Carregando extends Resultado {}
class Sucesso extends Resultado {
  final List<Leitura> dados;
  Sucesso(this.dados);
}
class Erro extends Resultado {
  final String mensagem;
  Erro(this.mensagem);
}

// switch exaustivo: sem default
String render(Resultado r) => switch (r) {
  Carregando() => 'Carregando...',
  Sucesso(:final dados) => '${dados.length} leituras',
  Erro(:final mensagem) => 'Falha: $mensagem',
};
\end{codigo}
    \end{column}
    \begin{column}{0.44\textwidth}
      \begin{conceito}[Por que \cod{sealed}?]
        Todos os subtipos ficam no mesmo arquivo; o compilador conhece a lista completa e exige que o \cod{switch} trate todos.
      \end{conceito}
      \begin{caixaambar}
        É exatamente o padrão usado para os estados de tela (carregando, sucesso, erro) nos Módulos 07 e 08.
      \end{caixaambar}
      \begin{dica}[Modificadores de classe]
        \cod{final}, \cod{base} e \cod{interface} controlam quem pode estender ou implementar a classe fora da biblioteca.
      \end{dica}
    \end{column}
  \end{columns}
\end{frame}

% =============================================================================
\section{Erros}

\begin{frame}[fragile]{Exceções}
  \begin{columns}[T]
    \begin{column}{0.52\textwidth}
\begin{codigo}[title=Tratando falhas]{Dart}
class SensorDesconectado implements Exception {
  final String id;
  SensorDesconectado(this.id);
  @override
  String toString() => 'Sensor $id desconectado';
}

double lerSensor(String id) {
  if (!conectado) throw SensorDesconectado(id);
  return 24.5;
}

try {
  final v = double.parse(entrada);
  print(lerSensor('t1') + v);
} on FormatException {
  print('valor inválido: $entrada');
} on SensorDesconectado catch (e) {
  print(e);
} finally {
  fecharArquivo();
}
\end{codigo}
    \end{column}
    \begin{column}{0.46\textwidth}
      \begin{conceito}[\cod{on} \emph{versus} \cod{catch}]
        \cod{on Tipo} trata um tipo específico; \cod{catch (e, pilha)} dá acesso ao objeto e à pilha de chamadas.
      \end{conceito}
      \begin{dica}
        \cod{double.tryParse(texto)} devolve \cod{null} em vez de lançar exceção. Ótimo para validar campos de formulário.
      \end{dica}
      \begin{atencao}[\cod{Exception} \emph{versus} \cod{Error}]
        \cod{Exception}: falhas esperadas (rede, formato) que devem ser tratadas. \cod{Error}: bugs de programação, que devem ser corrigidos.
      \end{atencao}
    \end{column}
  \end{columns}
\end{frame}

% =============================================================================
\section{Fechamento}

\begin{frame}{Resumo}
  \begin{resumo}
    \begin{itemize}
      \item Prefira \cod{final}; \cod{const} para valores de compilação; \cod{late} com cuidado.
      \item \cod{T?} marca valores anuláveis; use \cod{?.}, \cod{??} e promoção de tipo; evite \cod{!}.
      \item \cod{switch} de expressão com \emph{patterns} substitui cadeias de \cod{if}.
      \item Parâmetros nomeados e funções anônimas são a base da API do Flutter.
      \item Coleções com \cod{where}/\cod{map}/\cod{fold} e \emph{collection if/for}.
      \item Classes, \emph{mixins}, extensões, \emph{records} e \cod{sealed class} para modelar o domínio e os estados.
    \end{itemize}
  \end{resumo}
\end{frame}

\begin{frame}{Glossário}
  \begin{glossario}
    Paradigma & estilo de organizar o código (imperativo, objetos, funcional, declarativo) \\
    Tipagem estática & tipos verificados na compilação \\
    Solidez (\emph{soundness}) & garantia de que o tipo declarado vale também na execução \\
    \emph{Null safety} & distinção entre tipos anuláveis (\cod{String?}) e não anuláveis \\
    Promoção de tipo & o compilador refina o tipo após uma verificação (\cod{if (x != null)}) \\
    \emph{Mixin} & conjunto de membros reutilizado em várias classes sem herança \\
    \emph{Record} & valor anônimo e imutável com campos, comparado por valor \\
    \emph{Pattern} & forma que confere e desestrutura valores (\cod{switch}, \cod{if case}) \\
    Coleta de lixo & liberação automática de objetos sem referências \\
  \end{glossario}
\end{frame}

\begin{frame}{Questões de revisão}
  \begin{questoes}
    \item Que vantagens a tipagem estática e sólida do Dart traz para um app que vai crescer?
    \item Por que o nulo foi chamado de ``erro de um bilhão de dólares''? Como o Dart o trata?
    \item Qual a diferença entre \cod{final} e \cod{const}?
    \item Dê um exemplo de encapsulamento e um de polimorfismo no LabMonitor.
    \item O que é uma função pura? Por que ela é fácil de testar?
    \item Para que servem \cod{sealed class} e \cod{switch} exaustivo ao modelar o estado de uma tela?
    \item Qual a diferença entre igualdade de valor e de identidade? Como os \emph{records} se comportam?
    \item O que acontece com um objeto quando nenhuma variável aponta mais para ele?
  \end{questoes}
\end{frame}

\begin{frame}{Exercícios}
  \begin{exercicio}[Exercício 1 --- Conversões]
    Crie extensões \cod{cParaF}, \cod{fParaC} em \cod{double} e \cod{paraVolts} em \cod{int} (ADC de 12 bits, $V_{ref}$ = 3,3\,V).
  \end{exercicio}
  \begin{exercicio}[Exercício 2 --- Coleções]
    Dada uma lista de \cod{Leitura}, calcule a média por sensor (\cod{Map<String, double>}), as leituras fora de 10--40\,°C e a média móvel de 5 amostras do sensor \cod{'t1'}.
  \end{exercicio}
  \begin{exercicio}[Exercício 3 --- Estados]
    Modele com \cod{sealed class} os estados de uma conexão BLE (desconectado, procurando, conectando, conectado com RSSI, falha) e escreva um \cod{switch} exaustivo que gera a mensagem de cada estado.
  \end{exercicio}
  \begin{projeto}[LabMonitor --- Etapa 2]
    Crie \arq{lib/domain/models/} com \cod{Equipamento}, \cod{Bancada}, \cod{Leitura} e \cod{enum StatusBancada}, mais uma lista de dados fictícios.
  \end{projeto}
\end{frame}

\begin{frame}{Para aprofundar}
  \small
  \begin{itemize}
    \item Tour da linguagem: \url{https://dart.dev/language}
    \item Null safety: \url{https://dart.dev/null-safety}
    \item Patterns e records: \url{https://dart.dev/language/patterns}
    \item Modificadores de classe: \url{https://dart.dev/language/class-modifiers}
    \item Effective Dart (boas práticas): \url{https://dart.dev/effective-dart}
    \item DartPad: \url{https://dartpad.dev}
  \end{itemize}
\end{frame}

\begin{frame}{Respostas das questões de revisão}
  \begin{respostas}
    \item Erros de tipo aparecem na compilação, a IDE completa e refatora com segurança e, por ser sólida, a tipagem garante os tipos também em execução.
    \item Porque referências nulas causaram incontáveis falhas. O Dart separa \cod{String} de \cod{String?} e exige tratar o nulo (\cod{?.}, \cod{??}, promoção de tipo) antes do uso.
    \item \cod{final} é atribuída uma vez, em tempo de execução; \cod{const} é um valor conhecido na compilação, compartilhado e imutável.
    \item Encapsulamento: \cod{Reserva} só muda por \cod{cancelar()}. Polimorfismo: \cod{FonteBle} e \cod{FonteFake} implementam \cod{leituras()} cada uma do seu jeito.
    \item Devolve sempre o mesmo resultado para a mesma entrada, sem efeitos colaterais; o teste só confere entrada e saída.
    \item A \cod{sealed class} fecha o conjunto de estados (carregando, sucesso, erro) e o \cod{switch} exaustivo obriga a tratar todos.
    \item Identidade: o mesmo objeto (\cod{identical}); valor: o mesmo conteúdo. \emph{Records} com os mesmos campos são iguais por valor.
    \item Ele se torna inalcançável e é liberado pelo coletor de lixo em algum momento.
  \end{respostas}
\end{frame}

\encerramento{Módulo 03 (comum) --- Engenharia de software e Git: requisitos, métodos ágeis, versionamento e colaboração.}

\end{document}
